\documentclass[12pt]{article}
\usepackage[T2A]{fontenc}
\usepackage[utf8]{inputenc}
\usepackage[russian]{babel}
\usepackage{amsmath,amssymb}
\usepackage{geometry}
\geometry{margin=2.5cm}
\usepackage{indentfirst}

\title{К юбилею профессора Л.~Д.~Ландау:\\ полвека идей, опередивших время}
\author{А.~В.~Ковалёва, Д.~С.~Орлов, М.~П.~Соколова}
\date{2024}

\begin{document}
\maketitle

\section*{Аннотация}
Поздравительная статья, посвящённая юбилею выдающегося физика-теоретика. 

\vspace{0.5cm}
\noindent\textbf{Ключевые слова:} Ландау; юбилей; теоретическая физика; конденсированное состояние; история науки.

\vspace{0.3cm}
\noindent\textbf{Раздел журнала:} Поздравительный. \quad \textbf{Файл:} 1 из 2.

\section{Введение}
Идеи Л.~Д.~Ландау определили облик теоретической физики XX века. В этой статье мы вспоминаем ключевые результаты и их влияние на современную науку.

\section{Наследие}
Теория сверхтекучести, теория фазовых переходов второго рода и курс теоретической физики — лишь часть наследия, которое продолжает питать новые исследования.

\section{Заключение}
Мы поздравляем коллег и желаем новых открытий в духе Ландау.

\end{document}